\begin{abstract}
The Lightning Network routes payments across chains of untrusted
intermediaries using hashed time-locked contracts (HTLCs) layered on top of a
two-party channel revocation mechanism. Prior symbolic analyses have treated
these layers separately: two-party HTLC constructions on one side, generic
message-forwarding abstractions on the other. Their interaction --- revocation
and punishment together with multi-hop routing, per-hop fees, and staggered
block-height timing --- has not been mechanised as a single symbolic model.

We present such a model in \mtamarin, structured as five theories that together
discharge 87 lemma checks over 61 distinct properties under a full Dolev--Yao
adversary. Building it surfaced a modelling pitfall we believe generalises
beyond Lightning: encoding a one-shot on-chain output as a persistent fact
permits the corresponding rule to fire without bound, and every individual
safety lemma continues to hold while it does. We found two independent
instances of this in our own model --- unbounded punishment for a single act of
cheating, and unbounded HTLC reuse of one channel output --- diagnosed both as
the same defect, and repaired both with linear single-spend tokens. We state
the repair as an auditable discipline.

We further prove value conservation with per-hop fees, establish empirically
that the intermediary claim obligation is not implied by the timelock ordering
constraints, and reproduce the known wormhole attack. Because our model carries
amounts rather than abstracting them away, we pin the wormhole's theft to
exactly the sum of the bypassed fees --- a quantification the underlying
formalism must carry numbers to express. Finally, we lift the safety results to
paths of unbounded length by abstracting the routing layer, and give the
refinement argument that justifies the transfer.
\end{abstract}
